\section{A prospective longitudinal-gauge route}\label{app:research-route}
\subsection{Follow-up: fourth order inside the horizon}\label{sec:followup-late}
This is a proposed route, not an evaluated fourth-order asymptotic result. Its starting
point is the exact longitudinal-redshift identity already derived in the report.

\paragraph{Why change variables.} The sampled finite-wavelength maps have large cancelling
contributions in the temperature frame. No fourth-order decay law is established. The exploratory recombination in App.~\ref{app:numerics} suggests where
the residual cancellation occurs.
It sits between the field terms and the photon displacement, and it has two likely sources: the
static relabelling carried by $H_1\to-\frac6{k_0^2}(\partial\partial-\frac13\delta\nabla^2)g$,
and the $O(1)$ oscillation of the temperature-defined slicing, since the radiation
temperature itself oscillates inside the horizon. A purely spatial relabelling addresses the
first but not the second, and lowers the cancellation measure only from $10.3$ to $3.5$ at
$x_o=12$. In longitudinal variables both are removed at first order. The potentials
$\Phi=\Psi=-2(\sin x/x-\cos x)/x^2$ decay as $x^{-2}$, there is no first-order scalar shift,
and the $O(1)$ acoustic physics appears explicitly in the endpoint Sachs--Wolfe and Doppler
factors (Remark~\ref{rem:longitudinal}). The change is exact: the longitudinal redshift and
its identity with the temperature-frame exponent are already proved there.

\paragraph{What has to be built.}
\begin{enumerate}
\item The second- and third-order gauge generators (time shift and spatial vector) from the
  temperature frame to longitudinal gauge. At each order these solve linear equations with
  sources from lower orders, so the triangular scheme of \S\ref{sec:fields} applies. The
  existing field kernels supply the sources.
\item The longitudinal fields $\Phi_r$, $\Psi_r$, $\beta_r$ and $M_r$ through third order,
  obtained by transforming those kernels. Their late-time decay should be checked order by
  order: the induced second-order potentials are not guaranteed to decay as fast as the
  first-order ones.
\item The late-time expansion of the shell-integrated spectrum, with the non-uniform angles
  $\mu=\pm1$ (collinear) and $\mu=-\frac13$ (tensor sound cone) treated explicitly. Either
  use matched asymptotics in $(\mu,x_o)$, or subtract the resonant behaviour analytically
  before taking the limit. A pointwise limit followed by angular integration is not
  justified (see above).
\end{enumerate}

\paragraph{Checks available.} The result must agree with the temperature-frame map, which
is exact at finite times. It must reproduce the superhorizon coefficients of
\S\ref{sec:next-gradient} at small $x_o$, and its linear part must reproduce the first-order
subhorizon law of \S\ref{sec:subhorizon-first}. The second-order part can be
compared with the exploratory numerics, which reach $x_o=16$, as a consistency check. The
natural first milestone is the second-order part $\frac14C^{(22)}$ alone. It needs only the
second-order generator, and by Prop.~\ref{prop:nodes} it is the entire fourth-order answer at
the zeros of $T_\ell$.

